\section{泡沫、华人 AGI 与全球叙事}

前面都是公司和产品分析，最后一章处理外层变量：资本、泡沫和全球叙事。AGI 是否有泡沫，泡沫破裂的导火索是什么，人类和大猩猩智能差异在哪里，湾区新话题，以及“犹太人的金融，华人的 AGI”。这些话题看似分散，其实都在讨论 AI 的社会叙事和资本预期。

\begin{figure}[H]
\centering
\includegraphics[width=0.96\textwidth]{agi-bubble-dashboard.png}
\caption{AGI Bubble 仪表盘：泡沫是否破裂取决于收入、模型进步、CapEx 和产品兑现。自制概念图，依据 00:55:53--01:09:11 对谈内容整理。}
\end{figure}

\begin{knowledgebox}{读图：泡沫不只看情绪}
判断 AGI 泡沫，要看模型是否继续进步、收入是否兑现、CapEx 是否有回报、L4 产品是否继续出现、估值是否过度透支。情绪会波动，但基本面要看这些指标。
\end{knowledgebox}

\begin{knowledgebox}{泡沫触发器表}
\begin{tabularx}{\textwidth}{p{0.22\textwidth}p{0.34\textwidth}X}
\toprule
触发器 & 可能信号 & 为什么危险 \\
\midrule
模型进步放缓 & 新模型提升有限，成本更高 & 估值建立在持续能力跃迁上。 \\
收入不及预期 & 订阅、API、企业收入低于 CapEx & 资本开支难以被现金流解释。 \\
产品窗口关闭 & L4 体验迟迟不扩展 & 用户增长和付费意愿下降。 \\
硬件/能源约束 & GPU、数据中心、电力成本上升 & 训练和推理成本压缩利润。 \\
监管/版权风险 & 数据和生成结果被限制 & 数据和产品供给受影响。 \\
\bottomrule
\end{tabularx}
\end{knowledgebox}

\subsection{华人的 AGI}

“犹太人的金融，华人的 AGI”是一句带有叙事色彩的总结。它指向一个事实：全球 AGI 竞赛中，华人团队在模型、应用、工程和创业中都非常活跃。但挑战也明显：全球化品牌、资本市场、算力供应、企业销售和合规都不是纯技术问题。

\begin{figure}[H]
\centering
\includegraphics[width=0.94\textwidth]{chinese-agi-global.png}
\caption{华人的 AGI：全球 AGI 竞赛中，华人团队的技术、产品和商业能力被重新定价。自制概念图，依据 00:55:53--01:09:11 对谈内容整理。}
\end{figure}

\begin{warningbox}{叙事不能替代能力}
“华人的 AGI”可以是鼓舞性的叙事，但不能替代模型能力、产品能力、组织能力和全球商业能力。真正的全球公司需要同时解决技术、产品、资本、品牌和合规。
\end{warningbox}

\subsection{本章小结}

泡沫和全球叙事是 AI 周期的外层变量。模型进步和产品体验决定底层，资本和叙事决定波动，创业者必须同时看两层。

